% \iffalse meta-comment
%
% OTHR-CI -- LaTeX using the Corporate Design of OTH Regensburg
% ----------------------------------------------------------------------------
%
%  Copyright (C) 2015-2024 Michael Nietmetz <michael.niemetz@oth-regensburg.de>
%                 and the OTH Regensburg-LaTeX-Team
%                2025 Marei Peischl (peiTeX) <othr-ci@peitex.de>
%                 and Michael Nietmetz <michael.niemetz@oth-regensburg.de>
%
% ============================================================================
% This work may be distributed and/or modified under the
% conditions of the LaTeX Project Public License, either version 1.3c
% of this license or (at your option) any later version.
% The latest version of this license is in
% http://www.latex-project.org/lppl.txt
% and version 1.3c or later is part of all distributions of LaTeX
% version 2008/05/04 or later.
%
% This work has the LPPL maintenance status `maintained'.
%
% The Current Maintainers of this work are
%   Marei Peischl <othr-ci@peitex.de>
%   Michael Niemetz <michael.niemetz@oth-regensburg.de>
%
% The development repository can be found at
% https://codeberg.org/OTH-Regensburg/OTHRegensburg-LaTeX
% Please use the issue tracker for feedback!
% ============================================================================
%
% \fi
% \iffalse
%<*driver>
\ProvidesFile{othr-colors.dtx}
  [2025-05-22 v0.01-dev Color setup for  CI of OTH Regensburg]
%</driver>
%<@@=ptxcd>
%<*identification>
%<package>\NeedsTeXFormat{LaTeX2e}[2022-10-01]
%<package>\ProvidesExplPackage{othr-colors}{2025-05-22}{0.01-dev}{Color setup for  CI of OTH Regensburg}
%<definitions>\ProvidesFile{othr-colors.def}[2025-05-22 v0.01-dev Color definitions for OTHR-CI]
%</identification>
%<*driver>
\title{\textsf{othr-colors} --Color setup for the Corporate Design of OTH Regensburg%
  \thanks{This file describes version \fileversion, last revised \filedate.}
}
\providecommand*{\OTHRDocDTXfiles}{othr-colors.dtx}
% \iffalse meta-comment
%
% OTHR-CI -- LaTeX using the Corporate Design of OTH Regensburg
% ----------------------------------------------------------------------------
%
%  Copyright (C) 2015-2024 Michael Nietmetz <michael.niemetz@oth-regensburg.de>
%                 and the OTH Regensburg-LaTeX-Team
%                2025 Marei Peischl (peiTeX) <othr-ci@peitex.de>
%                 and Michael Nietmetz <michael.niemetz@oth-regensburg.de>
%
% ============================================================================
% This work may be distributed and/or modified under the
% conditions of the LaTeX Project Public License, either version 1.3c
% of this license or (at your option) any later version.
% The latest version of this license is in
% http://www.latex-project.org/lppl.txt
% and version 1.3c or later is part of all distributions of LaTeX
% version 2008/05/04 or later.
%
% This work has the LPPL maintenance status `maintained'.
%
% The Current Maintainers of this work are
%   Marei Peischl <othr-ci@peitex.de>
%   Michael Niemetz <michael.niemetz@oth-regensburg.de>
%
% The development repository can be found at
% https://codeberg.org/OTH-Regensburg/OTHRegensburg-LaTeX
% Please use the issue tracker for feedback!
% ============================================================================
%
% \fi
% \iffalse
%<*driver>
\typeout{Currently this file does not create documenation.}
%</driver>
%<@@=ptxcd>
%<*identification>
\NeedsTeXFormat{LaTeX2e}[2022-06-01]
%<driver>\ProvidesFile{othr-ci.dtx}
%<driver>  [2025-05-22 v0.01-dev  Fundamental configuration for OTHR-CI – LaTeX using the Coporate Design of OTH Regensburg]
%<package>\ProvidesExplPackage{OTHR}{2025-05-22}{0.01-dev}{Fundamental configuration for OTHR-CI – LaTeX using the Coporate Design of OTH Regensburg}
%</identification>
%<*driver>
\documentclass[lm-default=false,cs-break-nohyphen]{l3doc}

\usepackage{minted}
\usepackage{accsupp}
\setminted{breaksymbol={\BeginAccSupp{method=escape,ActualText={}}\tiny\ensuremath{\hookrightarrow}\EndAccSupp{}}}
\newminted[examplecode]{latex}{tabsize=3,gobble=1,escapeinside=||,breaklines}

\ExplSyntaxOn
\makeatletter
\tl_new:N \l__ptxtools_doc_values_tl
\newcommand*{\codefamily}{\MacroFont}
\DeclareTextFontCommand{\codefont}{\codefamily}
\providecommand{\option}[1]{\codefont{#1}}

\cs_new:Npn \__ptxctools_parse_key_option:w #1 #2 #3 #4 = #5 \q_stop {
  \begingroup
  \let\PrintDescribeOption\PrintDescribeKeyOption
  \tl_set:Nn \l__ptxtools_doc_values_tl {\small\IfBooleanTF{#3}{\textsf{#5}}{\textsf{(\codefont{#5})}}}
  \noindent\DescribeOption{#4}\hfill
  \small
  \tl_if_blank:nTF {#2} {\meta{\textrm{initially~unset}}} {
    \IfBooleanTF{#1}
    {#2}
    {(default:~\codefont{#2})}
  }\strut\par
  \endgroup
}

\newcommand*{\PrintDescribeKeyOption}[1]{
  {\MacroFont #1=}\rlap{\hskip\marginparsep\l__ptxtools_doc_values_tl}
}

\bool_new:N \l_ptxtools_last_was_describe_bool
\NewDocumentCommand{\DescribeKeyOption}{smms}{
  \bool_if:NTF \l_ptxtools_last_was_describe_bool {
    \vskip -\lastskip
  }{
    \par
    \dim_compare:nNnF {\parskip} > {\c_zero_dim} {\medskip}
  }
  \__ptxctools_parse_key_option:w #1 {#3} #4 #2 \q_stop

  \bool_set_true:N \l_ptxtools_last_was_describe_bool
  \AddToHookNext{para/begin}{
    \bool_set_false:N \l_ptxtools_last_was_describe_bool
    \OmitIndent
  }
  \smallskip
}

% to support small layout adjustments for the texdoc version
\newcommand*{\IfOnlyDocumentation}[1]{
  \bool_if_exist:NT \g__codedoc_typeset_implementation_bool {
    \bool_if:NF \g__codedoc_typeset_implementation_bool {#1}
  }
}

\NewDocumentCommand{\PTXToolsDocSection}{om}{
  \str_if_eq:NNF  \CurrentFile \OTHRDocDTXfiles {
    \section[#1]{#2}
  }
}
\ExplSyntaxOff

\usepackage{underscore-ltx}

\DeclareRobustCommand{\issue}[1]{%
  {\DeclareFieldFormat{citeurlpostfix}{\href{##1/issues/#1}{\##1}}\citefield{othr-ci-repo}[citeurlpostfix]{url}}%
}
\newcommand*{\version}[1]{v#1}

\EnableCrossrefs
% \CodelineIndex
\RecordChanges

\DoNotIndex{\newcommand,\newenvironment}

\expandafter\ifx\csname OTHRDocDTXfiles\endcsname\relax
  \title{OTHRegensburg-CI --- \LaTeX\ using the Corporate Design of OTH-Regensburg%
    \thanks{This file describes version \fileversion, last revised \filedate.}
  }
  \GetFileInfo{othr-ci.dtx}
\else
  \GetFileInfo{\OTHRDocDTXfiles}
% subtitle? \fileinfo?
\fi

\author{%
  Marei Peischl\thanks{Email: \href{mailto:othr-ci@peitex.de}{othr-ci@peitex.de}}%
  \and Michael Niemetz\thanks{Email: \href{mailto:michael.niemetz@oth-regensburg.de}{michael.niemetz@oth-regensburg.de}}%
}
\date{\fileversion~from \filedate}

\usepackage{OTHR}

\providecommand*{\OTHRDocDTXfiles}{
  othr-ci.dtx,
  othr-beamer.dtx,
  othr-colors.dtx,
  othr-fonts.dtx
}

\begin{document}

\maketitle
\EnableDocumentation\DisableImplementation

\begin{abstract}
  The OTHRegensburg-CI bundle provides possibilities to use the CI of OTH Regensburg for \LaTeX{}.

  Parts of this bundle are historically grown, but the beamer template was completely rewritten in 2025. This document is the starting point for a full documentation of the bundle, but all document types beside beamer are described within the old documentation, which is only available in German \cite{othrTeX}.
\end{abstract}

\tableofcontents

\iffalse
\section{Contained packages and supported document types}

This bundle supports the following document types:

\begin{description}
\item[Beamer slides], see section \ref{sec:beamer}
\end{description}
\fi

\UseName{exp_args:No} \DocInput {\OTHRDocDTXfiles}

\DisableDocumentation\EnableImplementation

\begin{implementation}
\section{Implementation}
\end{implementation}

\DocInputAgain

\PrintChanges
\PrintIndex
\end{document}
%</driver>
% \fi
% \iffalse
%<*declareoptions|options>
%<*class>
% \fi
% \begin{implementation}
%    \begin{macrocode}
\str_const:Nn \c_@@_base_str {
%<beamer>  beamer
}

\exp_args:Ne \keys_define:nn {ptxcd/\c__ptxcd_base_str} {
%    \end{macrocode}
% \iffalse
%</class>
%    \begin{macrocode}
%<wrapper>  \keys_define:nn {ptxcd/othr-ci} {
%<gobble>  }^^A gobble brace for latexindent
%    \end{macrocode}
% \fi
% \end{implementation}
% \begin{documentation}
% \section{General layout options}
% The OTHRegensburg-CI bundle supports multiple document types.
% While rewriting the beamer theme the options which are available within multiple classes or packages of this bundle have been constructe separately to simplfy maintenance.
%
% Currently these options are only fully supporte by the beamer class \cls{othr-beamer} but this might be extended in the future.
% \end{documentation}
% \begin{documentation}
% \DescribeKeyOption{department=\meta{departmentshorthand}}{default}
% As the CI guideline is attaching colors to specific departments as well as there is a mechanism to load department specific config files this option may be used in a lot of different ways.
% For example the \pkg{othr-colors} package is using it to set  a department specific highlighting color.%^^A TODO provide reference!!
% The \value{default} is used to determine if a user specific setting as provided or not.
% In the current setup it has no effect towards leaving it empty.
% \end{documentation}
% \begin{implementation}
% \begin{optionenv}{department}
%    \begin{macrocode}
  department .choice:,
  department / default .code:n = \str_gset:Nn \g_ptxcd_department_str {default},
%<class>  department .initial:n = default,
  department / unknown .code:n = {
      \str_gset:NV \g_ptxcd_department_str \l_keys_value_tl
    },
%    \end{macrocode}
% \end{optionenv}
% \end{implementation}
% \begin{documentation}
% \DescribeKeyOption{departmentconfigprefix=\meta{prefix}}{othr-}
% The loading mechanism to be used to load department specific config files allows to change the file prefix.
% This might be useful if a specific department wants some internal templates to be shared separately.
% In that case it is also possible to use this option to set a specific directory where those config files can be found.
% By default this setting will search for files called |othr-|\meta{department}|.cfg|, where |othr-| can be set using this option and \option{department} would configure the second part.
% \end{documentation}
% \begin{implementation}
% \begin{optionenv}{departmentconfigprefix}
%    \begin{macrocode}
  departmentconfigprefix .tl_gset:N = \g_@@_config_prefix_tl,
%<class>  departmentconfigprefix .initial:n = othr-,
%    \end{macrocode}
% \end{optionenv}
% \end{implementation}
% \begin{documentation}
% \DescribeKeyOption{logofile=\meta{<main-logo-filename/prefix>}}{othr-logo}
% \DescribeKeyOption{titlelogofile=\meta{filename}}{\meta{logofile}_flat}
% \DescribeKeyOption{logofile=\meta{filename}}{\meta{logofile}_short}
% The beamer template is using two different logo files.
% The \enquote{flat} variant is used for the title page and the \enquote{short} appears within the frame title.
% To simplify the global config the option \option{logofile} is used as a wrapper to configure both defaults.
%
% As the official logos are shipped within this bundle it should not be necessary for users to use these options.
%
% \end{documentation}
% \begin{implementation}
% \begin{optionenv}{logofile,short-logofile,flat-logofile}
%    \begin{macrocode}
  logofile .tl_gset:N = \g_ptxcd_logofile_tl,
%<class>  logofile .initial:n = othr-logo,
  short-logofile .tl_gset:N = \g_@@_logofile_short_tl,
%<class>  short-logofile .initial:n =  \g_ptxcd_logofile_tl _short,
  flat-logofile .tl_gset:N = \g_@@_logofile_flat_tl,
%<class>  flat-logofile .initial:n =  \g_ptxcd_logofile_tl _flat,
%    \end{macrocode}
% \end{optionenv}
% \end{implementation}
% \begin{documentation}
% \DescribeKeyOption{accept-missing-logos=\meta{boolean}}{false}
% This option exists for compatibility reasons.
% It allows to use the \pkg{othr-logos} package which has been existing the former logo variants.
% In case you need to compile old tex-files and get error messages about missing logos, this option can be enabled.
% \end{documentation}
% \begin{implementation}
% \begin{optionenv}{accept-missing-logos}
%    \begin{macrocode}
  accept-missing-logos .bool_gset:N = \g_ptxcd_logo_workaround_bool,
%<class>  accept-missing-logos .initial:n = false,
  accept-missing-logos .usage:n = load,
  accept-missing-logos .default:n = true,
%    \end{macrocode}
% \end{optionenv}
% \iffalse
%<*wrapper>
% \fi
% \begin{optionenv}{useDepartmentLogo,useOTHRLogo}
%   Compatibility options for logo package.
%    \begin{macrocode}
  useDepartmentLogo .code:n = \PassOptionsToPackage{useDepartmentLogo}{othr-logos},
  useOTHRLogo .code:n = \PassOptionsToPackage{useOTHRLogo}{othr-logos},
%    \end{macrocode}
% \end{optionenv}
% \iffalse
%</wrapper>
%</declareoptions|options>
% \fi
%    \begin{macrocode}
%<processoptions|init-department>}
%    \end{macrocode}
% \iffalse
%<*init-department>
% \fi
% Configure defaults, in case the packages are used without loading a class being part of  OTHR-CI.
%    \begin{macrocode}
\str_if_exist:NF \g_ptxcd_department_str {
  \str_new:N \g_ptxcd_department_str
  \str_gset:Nn \g_ptxcd_department_str {default}
  \tl_gset:Nn \g_ptxcd_logofile_tl {othr-logo}
  \tl_gset:Nn \g_@@_logofile_short_tl { \g_ptxcd_logofile_tl _short}
  \tl_gset:Nn \g_@@_logofile_flat_tl { \g_ptxcd_logofile_tl _flat}
  \tl_gset:Nn \g_@@_config_prefix_tl {othr-}
}
%    \end{macrocode}
% \iffalse
%</init-department>
%<*processoptions>
% \fi
%    \begin{macrocode}
%<!wrapper>\ProcessKeyOptions[ptxcd/\c__ptxcd_base_str]
%<wrapper>\ProcessKeyOptions[ptxcd/othr-ci]
%    \end{macrocode}
% \iffalse
%</processoptions>
%<*wrapper&main>
% \fi
%    \begin{macrocode}
\RequirePackage{othr-colors}
\RequirePackage{othr-logos}
\RequirePackage{othr-fonts}
%    \end{macrocode}
% \iffalse
%</wrapper&main>
% \fi
% \end{implementation}
\endinput

%</driver>
% \fi
% \PTXToolsDocSection{Colors}
% \begin{documentation}
%
% The colors can be used like within standard \LaTeX.
% The predefined OTHR-CI specific colors are:
% \begin{multicols}{4}
%   \raggedright
%   \ExplSyntaxOn
%   \RenewDocumentCommand{\ProvidesFile}{mo}{}
%   \RenewDocumentCommand{\definecolor}{ommm}{\makebox[5em]{#2}\textcolor{#2}{\rule{1cm}{\ht\strutbox}}\par\ignorespaces}
%   \let\endinput\relax
%   \file_get:nnN {othr-colors.def} {\cctab_select:N \c_code_cctab} \l_tmpa_tl
%   \l_tmpa_tl
%   \ExplSyntaxOff
% \end{multicols}
% \end{documentation}
% \begin{implementation}
% \iffalse
%<*declareoptions>
% \fi
%    \begin{macrocode}
\str_if_exist:NF \g_ptxcd_department_str {
  \str_new:N \g_ptxcd_department_str
  \str_gset:Nn \g_ptxcd_department_str {default}
}

\bool_new:N \g_@@_color_blacktext_bool
\keys_define:nn {ptxcd/colors} {
%    \end{macrocode}
% \iffalse
%</declareoptions>
%<*options>
% \fi
% \end{implementation}
%
% \begin{documentation}
% \DescribeKeyOption{accentcolor=\meta{Color}}{0b}*
% \DescribeOption{accent=}
%
% Highlight color used for all highlight elements by default.
% The options below will reference this color as  default setting meta{accentcolor}.
% Depending on the documentclass it might be used for the so-called \enquote{identbar} as well as highlighted test elements, e.\,g. within letters if \option{premium} is enabled.
%
% |accent| is an alias.
% We recommend using |accentcolor| to simplify recognizing the purpose of that option.
% \end{documentation}
%
% \begin{implementation}
% \begin{optionenv}{accent=}
%    \begin{macrocode}
  accent .tl_gset:N = \g_ptxcd_color_accent_tl,
%<package>  accent .initial:n = 0b,
  accentcolor .meta:n = {accent =#1},
%    \end{macrocode}
% \end{optionenv}^^A accent
% \end{implementation}
%
% \begin{documentation}
% \DescribeKeyOption{textaccent=\meta{color}}{\meta{accentcolor}}*
% The accent color for highlighted text. Usage depends on the document class.
% \end{documentation}
%
% \begin{implementation}
% \begin{optionenv}{textaccent=\meta{color}}
%    \begin{macrocode}
  textaccent .tl_gset:N = \g_ptxcd_color_textaccent_tl,
%<package>  textaccent .initial:n = \g_ptxcd_color_accent_tl,
  textaccentcolor .meta:n = {textaccent = #1},
%    \end{macrocode}
% \end{optionenv}^^A textaccent
% \end{implementation}
%
% \begin{documentation}
% \DescribeKeyOption{text=black/default}{default}*
% The text placed on top a colored area, e.\,g. colored title background, is defined to have enough contrast.
% \end{documentation}
%
% \begin{implementation}
% \begin{optionenv}{text}
%   \iffalse
%<*package>
% \fi
%    \begin{macrocode}
  text .choice:,
  text/black .code:n = \bool_gset_true:N \g_@@_color_blacktext_bool,
  text/default .code:n = \bool_gset_false:N \g_@@_color_blacktext_bool,
  text .initial:n = {default},
%    \end{macrocode}
%   \iffalse
%</package>
% \fi
%    \begin{macrocode}
%<!package>  text .code:n = \exp_args:Ne \PassOptionsToPackage{text=\l_keys_value_tl}{othr-colors},
%    \end{macrocode}
%   For compatibility reasons there is an alias for \option{text=black} called \option{blackFont}.
%    \begin{macrocode}
  blackFont.meta:n = {text=black},
  blackFont .value_forbidden:n = true,
%    \end{macrocode}
% \end{optionenv}^^A text
% \end{implementation}
% \begin{documentation}
% \DescribeKeyOption*{colormode=cmyk/RGB}{Depends on document type}*
% The colormode can be selected to overwrite the default setting.
% Please be aware, that there is a difference in casing and \code{RGB} has to be used uppercase and \code{cmyk} lowercase.
% The reason for this is that internally uppercase and lowercase model names slightly differ because of rounding.
% \end{documentation}
% \begin{implementation}
% \begin{optionenv}{colormode}
%    \begin{macrocode}
  colormode .choices:nn = {cmyk,RGB}{
      \PassOptionsToPackage{\l_keys_value_tl}{xcolor}
    },
  colormode .usage:n = load,
  RGB .meta:n = {colormode=RGB},
  cymk .meta:n = {colormode=cmyk},
%    \end{macrocode}
% \end{optionenv}
% \iffalse
%<*package>
% \fi
%    \begin{macrocode}
}
\ProcessKeyOptions[ptxcd/colors]
\RequirePackage{xcolor}

% \iffalse meta-comment
%
% OTHR-CI -- LaTeX using the Corporate Design of OTH Regensburg
% ----------------------------------------------------------------------------
%
%  Copyright (C) 2015-2024 Michael Nietmetz <michael.niemetz@oth-regensburg.de>
%                 and the OTH Regensburg-LaTeX-Team
%                2025 Marei Peischl (peiTeX) <othr-ci@peitex.de>
%                 and Michael Nietmetz <michael.niemetz@oth-regensburg.de>
%
% ============================================================================
% This work may be distributed and/or modified under the
% conditions of the LaTeX Project Public License, either version 1.3c
% of this license or (at your option) any later version.
% The latest version of this license is in
% http://www.latex-project.org/lppl.txt
% and version 1.3c or later is part of all distributions of LaTeX
% version 2008/05/04 or later.
%
% This work has the LPPL maintenance status `maintained'.
%
% The Current Maintainers of this work are
%   Marei Peischl <othr-ci@peitex.de>
%   Michael Niemetz <michael.niemetz@oth-regensburg.de>
%
% The development repository can be found at
% https://codeberg.org/OTH-Regensburg/OTHRegensburg-LaTeX
% Please use the issue tracker for feedback!
% ============================================================================
%
% \fi
% \iffalse
%<*driver>
\ProvidesFile{othr-colors.dtx}
  [2025-05-22 v0.01-dev Color setup for  CI of OTH Regensburg]
%</driver>
%<@@=ptxcd>
%<*identification>
%<package>\NeedsTeXFormat{LaTeX2e}[2022-10-01]
%<package>\ProvidesExplPackage{othr-colors}{2025-05-22}{0.01-dev}{Color setup for  CI of OTH Regensburg}
%<definitions>\ProvidesFile{othr-colors.def}[2025-05-22 v0.01-dev Color definitions for OTHR-CI]
%</identification>
%<*driver>
\title{\textsf{othr-colors} --Color setup for the Corporate Design of OTH Regensburg%
  \thanks{This file describes version \fileversion, last revised \filedate.}
}
\providecommand*{\OTHRDocDTXfiles}{othr-colors.dtx}
% \iffalse meta-comment
%
% OTHR-CI -- LaTeX using the Corporate Design of OTH Regensburg
% ----------------------------------------------------------------------------
%
%  Copyright (C) 2015-2024 Michael Nietmetz <michael.niemetz@oth-regensburg.de>
%                 and the OTH Regensburg-LaTeX-Team
%                2025 Marei Peischl (peiTeX) <othr-ci@peitex.de>
%                 and Michael Nietmetz <michael.niemetz@oth-regensburg.de>
%
% ============================================================================
% This work may be distributed and/or modified under the
% conditions of the LaTeX Project Public License, either version 1.3c
% of this license or (at your option) any later version.
% The latest version of this license is in
% http://www.latex-project.org/lppl.txt
% and version 1.3c or later is part of all distributions of LaTeX
% version 2008/05/04 or later.
%
% This work has the LPPL maintenance status `maintained'.
%
% The Current Maintainers of this work are
%   Marei Peischl <othr-ci@peitex.de>
%   Michael Niemetz <michael.niemetz@oth-regensburg.de>
%
% The development repository can be found at
% https://codeberg.org/OTH-Regensburg/OTHRegensburg-LaTeX
% Please use the issue tracker for feedback!
% ============================================================================
%
% \fi
% \iffalse
%<*driver>
\typeout{Currently this file does not create documenation.}
%</driver>
%<@@=ptxcd>
%<*identification>
\NeedsTeXFormat{LaTeX2e}[2022-06-01]
%<driver>\ProvidesFile{othr-ci.dtx}
%<driver>  [2025-05-22 v0.01-dev  Fundamental configuration for OTHR-CI – LaTeX using the Coporate Design of OTH Regensburg]
%<package>\ProvidesExplPackage{OTHR}{2025-05-22}{0.01-dev}{Fundamental configuration for OTHR-CI – LaTeX using the Coporate Design of OTH Regensburg}
%</identification>
%<*driver>
\documentclass[lm-default=false,cs-break-nohyphen]{l3doc}

\usepackage{minted}
\usepackage{accsupp}
\setminted{breaksymbol={\BeginAccSupp{method=escape,ActualText={}}\tiny\ensuremath{\hookrightarrow}\EndAccSupp{}}}
\newminted[examplecode]{latex}{tabsize=3,gobble=1,escapeinside=||,breaklines}

\ExplSyntaxOn
\makeatletter
\tl_new:N \l__ptxtools_doc_values_tl
\newcommand*{\codefamily}{\MacroFont}
\DeclareTextFontCommand{\codefont}{\codefamily}
\providecommand{\option}[1]{\codefont{#1}}

\cs_new:Npn \__ptxctools_parse_key_option:w #1 #2 #3 #4 = #5 \q_stop {
  \begingroup
  \let\PrintDescribeOption\PrintDescribeKeyOption
  \tl_set:Nn \l__ptxtools_doc_values_tl {\small\IfBooleanTF{#3}{\textsf{#5}}{\textsf{(\codefont{#5})}}}
  \noindent\DescribeOption{#4}\hfill
  \small
  \tl_if_blank:nTF {#2} {\meta{\textrm{initially~unset}}} {
    \IfBooleanTF{#1}
    {#2}
    {(default:~\codefont{#2})}
  }\strut\par
  \endgroup
}

\newcommand*{\PrintDescribeKeyOption}[1]{
  {\MacroFont #1=}\rlap{\hskip\marginparsep\l__ptxtools_doc_values_tl}
}

\bool_new:N \l_ptxtools_last_was_describe_bool
\NewDocumentCommand{\DescribeKeyOption}{smms}{
  \bool_if:NTF \l_ptxtools_last_was_describe_bool {
    \vskip -\lastskip
  }{
    \par
    \dim_compare:nNnF {\parskip} > {\c_zero_dim} {\medskip}
  }
  \__ptxctools_parse_key_option:w #1 {#3} #4 #2 \q_stop

  \bool_set_true:N \l_ptxtools_last_was_describe_bool
  \AddToHookNext{para/begin}{
    \bool_set_false:N \l_ptxtools_last_was_describe_bool
    \OmitIndent
  }
  \smallskip
}

% to support small layout adjustments for the texdoc version
\newcommand*{\IfOnlyDocumentation}[1]{
  \bool_if_exist:NT \g__codedoc_typeset_implementation_bool {
    \bool_if:NF \g__codedoc_typeset_implementation_bool {#1}
  }
}

\NewDocumentCommand{\PTXToolsDocSection}{om}{
  \str_if_eq:NNF  \CurrentFile \OTHRDocDTXfiles {
    \section[#1]{#2}
  }
}
\ExplSyntaxOff

\usepackage{underscore-ltx}

\DeclareRobustCommand{\issue}[1]{%
  {\DeclareFieldFormat{citeurlpostfix}{\href{##1/issues/#1}{\##1}}\citefield{othr-ci-repo}[citeurlpostfix]{url}}%
}
\newcommand*{\version}[1]{v#1}

\EnableCrossrefs
% \CodelineIndex
\RecordChanges

\DoNotIndex{\newcommand,\newenvironment}

\expandafter\ifx\csname OTHRDocDTXfiles\endcsname\relax
  \title{OTHRegensburg-CI --- \LaTeX\ using the Corporate Design of OTH-Regensburg%
    \thanks{This file describes version \fileversion, last revised \filedate.}
  }
  \GetFileInfo{othr-ci.dtx}
\else
  \GetFileInfo{\OTHRDocDTXfiles}
% subtitle? \fileinfo?
\fi

\author{%
  Marei Peischl\thanks{Email: \href{mailto:othr-ci@peitex.de}{othr-ci@peitex.de}}%
  \and Michael Niemetz\thanks{Email: \href{mailto:michael.niemetz@oth-regensburg.de}{michael.niemetz@oth-regensburg.de}}%
}
\date{\fileversion~from \filedate}

\usepackage{OTHR}

\providecommand*{\OTHRDocDTXfiles}{
  othr-ci.dtx,
  othr-beamer.dtx,
  othr-colors.dtx,
  othr-fonts.dtx
}

\begin{document}

\maketitle
\EnableDocumentation\DisableImplementation

\begin{abstract}
  The OTHRegensburg-CI bundle provides possibilities to use the CI of OTH Regensburg for \LaTeX{}.

  Parts of this bundle are historically grown, but the beamer template was completely rewritten in 2025. This document is the starting point for a full documentation of the bundle, but all document types beside beamer are described within the old documentation, which is only available in German \cite{othrTeX}.
\end{abstract}

\tableofcontents

\iffalse
\section{Contained packages and supported document types}

This bundle supports the following document types:

\begin{description}
\item[Beamer slides], see section \ref{sec:beamer}
\end{description}
\fi

\UseName{exp_args:No} \DocInput {\OTHRDocDTXfiles}

\DisableDocumentation\EnableImplementation

\begin{implementation}
\section{Implementation}
\end{implementation}

\DocInputAgain

\PrintChanges
\PrintIndex
\end{document}
%</driver>
% \fi
% \iffalse
%<*declareoptions|options>
%<*class>
% \fi
% \begin{implementation}
%    \begin{macrocode}
\str_const:Nn \c_@@_base_str {
%<beamer>  beamer
}

\exp_args:Ne \keys_define:nn {ptxcd/\c__ptxcd_base_str} {
%    \end{macrocode}
% \iffalse
%</class>
%    \begin{macrocode}
%<wrapper>  \keys_define:nn {ptxcd/othr-ci} {
%<gobble>  }^^A gobble brace for latexindent
%    \end{macrocode}
% \fi
% \end{implementation}
% \begin{documentation}
% \section{General layout options}
% The OTHRegensburg-CI bundle supports multiple document types.
% While rewriting the beamer theme the options which are available within multiple classes or packages of this bundle have been constructe separately to simplfy maintenance.
%
% Currently these options are only fully supporte by the beamer class \cls{othr-beamer} but this might be extended in the future.
% \end{documentation}
% \begin{documentation}
% \DescribeKeyOption{department=\meta{departmentshorthand}}{default}
% As the CI guideline is attaching colors to specific departments as well as there is a mechanism to load department specific config files this option may be used in a lot of different ways.
% For example the \pkg{othr-colors} package is using it to set  a department specific highlighting color.%^^A TODO provide reference!!
% The \value{default} is used to determine if a user specific setting as provided or not.
% In the current setup it has no effect towards leaving it empty.
% \end{documentation}
% \begin{implementation}
% \begin{optionenv}{department}
%    \begin{macrocode}
  department .choice:,
  department / default .code:n = \str_gset:Nn \g_ptxcd_department_str {default},
%<class>  department .initial:n = default,
  department / unknown .code:n = {
      \str_gset:NV \g_ptxcd_department_str \l_keys_value_tl
    },
%    \end{macrocode}
% \end{optionenv}
% \end{implementation}
% \begin{documentation}
% \DescribeKeyOption{departmentconfigprefix=\meta{prefix}}{othr-}
% The loading mechanism to be used to load department specific config files allows to change the file prefix.
% This might be useful if a specific department wants some internal templates to be shared separately.
% In that case it is also possible to use this option to set a specific directory where those config files can be found.
% By default this setting will search for files called |othr-|\meta{department}|.cfg|, where |othr-| can be set using this option and \option{department} would configure the second part.
% \end{documentation}
% \begin{implementation}
% \begin{optionenv}{departmentconfigprefix}
%    \begin{macrocode}
  departmentconfigprefix .tl_gset:N = \g_@@_config_prefix_tl,
%<class>  departmentconfigprefix .initial:n = othr-,
%    \end{macrocode}
% \end{optionenv}
% \end{implementation}
% \begin{documentation}
% \DescribeKeyOption{logofile=\meta{<main-logo-filename/prefix>}}{othr-logo}
% \DescribeKeyOption{titlelogofile=\meta{filename}}{\meta{logofile}_flat}
% \DescribeKeyOption{logofile=\meta{filename}}{\meta{logofile}_short}
% The beamer template is using two different logo files.
% The \enquote{flat} variant is used for the title page and the \enquote{short} appears within the frame title.
% To simplify the global config the option \option{logofile} is used as a wrapper to configure both defaults.
%
% As the official logos are shipped within this bundle it should not be necessary for users to use these options.
%
% \end{documentation}
% \begin{implementation}
% \begin{optionenv}{logofile,short-logofile,flat-logofile}
%    \begin{macrocode}
  logofile .tl_gset:N = \g_ptxcd_logofile_tl,
%<class>  logofile .initial:n = othr-logo,
  short-logofile .tl_gset:N = \g_@@_logofile_short_tl,
%<class>  short-logofile .initial:n =  \g_ptxcd_logofile_tl _short,
  flat-logofile .tl_gset:N = \g_@@_logofile_flat_tl,
%<class>  flat-logofile .initial:n =  \g_ptxcd_logofile_tl _flat,
%    \end{macrocode}
% \end{optionenv}
% \end{implementation}
% \begin{documentation}
% \DescribeKeyOption{accept-missing-logos=\meta{boolean}}{false}
% This option exists for compatibility reasons.
% It allows to use the \pkg{othr-logos} package which has been existing the former logo variants.
% In case you need to compile old tex-files and get error messages about missing logos, this option can be enabled.
% \end{documentation}
% \begin{implementation}
% \begin{optionenv}{accept-missing-logos}
%    \begin{macrocode}
  accept-missing-logos .bool_gset:N = \g_ptxcd_logo_workaround_bool,
%<class>  accept-missing-logos .initial:n = false,
  accept-missing-logos .usage:n = load,
  accept-missing-logos .default:n = true,
%    \end{macrocode}
% \end{optionenv}
% \iffalse
%<*wrapper>
% \fi
% \begin{optionenv}{useDepartmentLogo,useOTHRLogo}
%   Compatibility options for logo package.
%    \begin{macrocode}
  useDepartmentLogo .code:n = \PassOptionsToPackage{useDepartmentLogo}{othr-logos},
  useOTHRLogo .code:n = \PassOptionsToPackage{useOTHRLogo}{othr-logos},
%    \end{macrocode}
% \end{optionenv}
% \iffalse
%</wrapper>
%</declareoptions|options>
% \fi
%    \begin{macrocode}
%<processoptions|init-department>}
%    \end{macrocode}
% \iffalse
%<*init-department>
% \fi
% Configure defaults, in case the packages are used without loading a class being part of  OTHR-CI.
%    \begin{macrocode}
\str_if_exist:NF \g_ptxcd_department_str {
  \str_new:N \g_ptxcd_department_str
  \str_gset:Nn \g_ptxcd_department_str {default}
  \tl_gset:Nn \g_ptxcd_logofile_tl {othr-logo}
  \tl_gset:Nn \g_@@_logofile_short_tl { \g_ptxcd_logofile_tl _short}
  \tl_gset:Nn \g_@@_logofile_flat_tl { \g_ptxcd_logofile_tl _flat}
  \tl_gset:Nn \g_@@_config_prefix_tl {othr-}
}
%    \end{macrocode}
% \iffalse
%</init-department>
%<*processoptions>
% \fi
%    \begin{macrocode}
%<!wrapper>\ProcessKeyOptions[ptxcd/\c__ptxcd_base_str]
%<wrapper>\ProcessKeyOptions[ptxcd/othr-ci]
%    \end{macrocode}
% \iffalse
%</processoptions>
%<*wrapper&main>
% \fi
%    \begin{macrocode}
\RequirePackage{othr-colors}
\RequirePackage{othr-logos}
\RequirePackage{othr-fonts}
%    \end{macrocode}
% \iffalse
%</wrapper&main>
% \fi
% \end{implementation}
\endinput

%</driver>
% \fi
% \PTXToolsDocSection{Colors}
% \begin{documentation}
%
% The colors can be used like within standard \LaTeX.
% The predefined OTHR-CI specific colors are:
% \begin{multicols}{4}
%   \raggedright
%   \ExplSyntaxOn
%   \RenewDocumentCommand{\ProvidesFile}{mo}{}
%   \RenewDocumentCommand{\definecolor}{ommm}{\makebox[5em]{#2}\textcolor{#2}{\rule{1cm}{\ht\strutbox}}\par\ignorespaces}
%   \let\endinput\relax
%   \file_get:nnN {othr-colors.def} {\cctab_select:N \c_code_cctab} \l_tmpa_tl
%   \l_tmpa_tl
%   \ExplSyntaxOff
% \end{multicols}
% \end{documentation}
% \begin{implementation}
% \iffalse
%<*declareoptions>
% \fi
%    \begin{macrocode}
\str_if_exist:NF \g_ptxcd_department_str {
  \str_new:N \g_ptxcd_department_str
  \str_gset:Nn \g_ptxcd_department_str {default}
}

\bool_new:N \g_@@_color_blacktext_bool
\keys_define:nn {ptxcd/colors} {
%    \end{macrocode}
% \iffalse
%</declareoptions>
%<*options>
% \fi
% \end{implementation}
%
% \begin{documentation}
% \DescribeKeyOption{accentcolor=\meta{Color}}{0b}*
% \DescribeOption{accent=}
%
% Highlight color used for all highlight elements by default.
% The options below will reference this color as  default setting meta{accentcolor}.
% Depending on the documentclass it might be used for the so-called \enquote{identbar} as well as highlighted test elements, e.\,g. within letters if \option{premium} is enabled.
%
% |accent| is an alias.
% We recommend using |accentcolor| to simplify recognizing the purpose of that option.
% \end{documentation}
%
% \begin{implementation}
% \begin{optionenv}{accent=}
%    \begin{macrocode}
  accent .tl_gset:N = \g_ptxcd_color_accent_tl,
%<package>  accent .initial:n = 0b,
  accentcolor .meta:n = {accent =#1},
%    \end{macrocode}
% \end{optionenv}^^A accent
% \end{implementation}
%
% \begin{documentation}
% \DescribeKeyOption{textaccent=\meta{color}}{\meta{accentcolor}}*
% The accent color for highlighted text. Usage depends on the document class.
% \end{documentation}
%
% \begin{implementation}
% \begin{optionenv}{textaccent=\meta{color}}
%    \begin{macrocode}
  textaccent .tl_gset:N = \g_ptxcd_color_textaccent_tl,
%<package>  textaccent .initial:n = \g_ptxcd_color_accent_tl,
  textaccentcolor .meta:n = {textaccent = #1},
%    \end{macrocode}
% \end{optionenv}^^A textaccent
% \end{implementation}
%
% \begin{documentation}
% \DescribeKeyOption{text=black/default}{default}*
% The text placed on top a colored area, e.\,g. colored title background, is defined to have enough contrast.
% \end{documentation}
%
% \begin{implementation}
% \begin{optionenv}{text}
%   \iffalse
%<*package>
% \fi
%    \begin{macrocode}
  text .choice:,
  text/black .code:n = \bool_gset_true:N \g_@@_color_blacktext_bool,
  text/default .code:n = \bool_gset_false:N \g_@@_color_blacktext_bool,
  text .initial:n = {default},
%    \end{macrocode}
%   \iffalse
%</package>
% \fi
%    \begin{macrocode}
%<!package>  text .code:n = \exp_args:Ne \PassOptionsToPackage{text=\l_keys_value_tl}{othr-colors},
%    \end{macrocode}
%   For compatibility reasons there is an alias for \option{text=black} called \option{blackFont}.
%    \begin{macrocode}
  blackFont.meta:n = {text=black},
  blackFont .value_forbidden:n = true,
%    \end{macrocode}
% \end{optionenv}^^A text
% \end{implementation}
% \begin{documentation}
% \DescribeKeyOption*{colormode=cmyk/RGB}{Depends on document type}*
% The colormode can be selected to overwrite the default setting.
% Please be aware, that there is a difference in casing and \code{RGB} has to be used uppercase and \code{cmyk} lowercase.
% The reason for this is that internally uppercase and lowercase model names slightly differ because of rounding.
% \end{documentation}
% \begin{implementation}
% \begin{optionenv}{colormode}
%    \begin{macrocode}
  colormode .choices:nn = {cmyk,RGB}{
      \PassOptionsToPackage{\l_keys_value_tl}{xcolor}
    },
  colormode .usage:n = load,
  RGB .meta:n = {colormode=RGB},
  cymk .meta:n = {colormode=cmyk},
%    \end{macrocode}
% \end{optionenv}
% \iffalse
%<*package>
% \fi
%    \begin{macrocode}
}
\ProcessKeyOptions[ptxcd/colors]
\RequirePackage{xcolor}

% \iffalse meta-comment
%
% OTHR-CI -- LaTeX using the Corporate Design of OTH Regensburg
% ----------------------------------------------------------------------------
%
%  Copyright (C) 2015-2024 Michael Nietmetz <michael.niemetz@oth-regensburg.de>
%                 and the OTH Regensburg-LaTeX-Team
%                2025 Marei Peischl (peiTeX) <othr-ci@peitex.de>
%                 and Michael Nietmetz <michael.niemetz@oth-regensburg.de>
%
% ============================================================================
% This work may be distributed and/or modified under the
% conditions of the LaTeX Project Public License, either version 1.3c
% of this license or (at your option) any later version.
% The latest version of this license is in
% http://www.latex-project.org/lppl.txt
% and version 1.3c or later is part of all distributions of LaTeX
% version 2008/05/04 or later.
%
% This work has the LPPL maintenance status `maintained'.
%
% The Current Maintainers of this work are
%   Marei Peischl <othr-ci@peitex.de>
%   Michael Niemetz <michael.niemetz@oth-regensburg.de>
%
% The development repository can be found at
% https://codeberg.org/OTH-Regensburg/OTHRegensburg-LaTeX
% Please use the issue tracker for feedback!
% ============================================================================
%
% \fi
% \iffalse
%<*driver>
\ProvidesFile{othr-colors.dtx}
  [2025-05-22 v0.01-dev Color setup for  CI of OTH Regensburg]
%</driver>
%<@@=ptxcd>
%<*identification>
%<package>\NeedsTeXFormat{LaTeX2e}[2022-10-01]
%<package>\ProvidesExplPackage{othr-colors}{2025-05-22}{0.01-dev}{Color setup for  CI of OTH Regensburg}
%<definitions>\ProvidesFile{othr-colors.def}[2025-05-22 v0.01-dev Color definitions for OTHR-CI]
%</identification>
%<*driver>
\title{\textsf{othr-colors} --Color setup for the Corporate Design of OTH Regensburg%
  \thanks{This file describes version \fileversion, last revised \filedate.}
}
\providecommand*{\OTHRDocDTXfiles}{othr-colors.dtx}
% \iffalse meta-comment
%
% OTHR-CI -- LaTeX using the Corporate Design of OTH Regensburg
% ----------------------------------------------------------------------------
%
%  Copyright (C) 2015-2024 Michael Nietmetz <michael.niemetz@oth-regensburg.de>
%                 and the OTH Regensburg-LaTeX-Team
%                2025 Marei Peischl (peiTeX) <othr-ci@peitex.de>
%                 and Michael Nietmetz <michael.niemetz@oth-regensburg.de>
%
% ============================================================================
% This work may be distributed and/or modified under the
% conditions of the LaTeX Project Public License, either version 1.3c
% of this license or (at your option) any later version.
% The latest version of this license is in
% http://www.latex-project.org/lppl.txt
% and version 1.3c or later is part of all distributions of LaTeX
% version 2008/05/04 or later.
%
% This work has the LPPL maintenance status `maintained'.
%
% The Current Maintainers of this work are
%   Marei Peischl <othr-ci@peitex.de>
%   Michael Niemetz <michael.niemetz@oth-regensburg.de>
%
% The development repository can be found at
% https://codeberg.org/OTH-Regensburg/OTHRegensburg-LaTeX
% Please use the issue tracker for feedback!
% ============================================================================
%
% \fi
% \iffalse
%<*driver>
\typeout{Currently this file does not create documenation.}
%</driver>
%<@@=ptxcd>
%<*identification>
\NeedsTeXFormat{LaTeX2e}[2022-06-01]
%<driver>\ProvidesFile{othr-ci.dtx}
%<driver>  [2025-05-22 v0.01-dev  Fundamental configuration for OTHR-CI – LaTeX using the Coporate Design of OTH Regensburg]
%<package>\ProvidesExplPackage{OTHR}{2025-05-22}{0.01-dev}{Fundamental configuration for OTHR-CI – LaTeX using the Coporate Design of OTH Regensburg}
%</identification>
%<*driver>
\documentclass[lm-default=false,cs-break-nohyphen]{l3doc}

\usepackage{minted}
\usepackage{accsupp}
\setminted{breaksymbol={\BeginAccSupp{method=escape,ActualText={}}\tiny\ensuremath{\hookrightarrow}\EndAccSupp{}}}
\newminted[examplecode]{latex}{tabsize=3,gobble=1,escapeinside=||,breaklines}

\ExplSyntaxOn
\makeatletter
\tl_new:N \l__ptxtools_doc_values_tl
\newcommand*{\codefamily}{\MacroFont}
\DeclareTextFontCommand{\codefont}{\codefamily}
\providecommand{\option}[1]{\codefont{#1}}

\cs_new:Npn \__ptxctools_parse_key_option:w #1 #2 #3 #4 = #5 \q_stop {
  \begingroup
  \let\PrintDescribeOption\PrintDescribeKeyOption
  \tl_set:Nn \l__ptxtools_doc_values_tl {\small\IfBooleanTF{#3}{\textsf{#5}}{\textsf{(\codefont{#5})}}}
  \noindent\DescribeOption{#4}\hfill
  \small
  \tl_if_blank:nTF {#2} {\meta{\textrm{initially~unset}}} {
    \IfBooleanTF{#1}
    {#2}
    {(default:~\codefont{#2})}
  }\strut\par
  \endgroup
}

\newcommand*{\PrintDescribeKeyOption}[1]{
  {\MacroFont #1=}\rlap{\hskip\marginparsep\l__ptxtools_doc_values_tl}
}

\bool_new:N \l_ptxtools_last_was_describe_bool
\NewDocumentCommand{\DescribeKeyOption}{smms}{
  \bool_if:NTF \l_ptxtools_last_was_describe_bool {
    \vskip -\lastskip
  }{
    \par
    \dim_compare:nNnF {\parskip} > {\c_zero_dim} {\medskip}
  }
  \__ptxctools_parse_key_option:w #1 {#3} #4 #2 \q_stop

  \bool_set_true:N \l_ptxtools_last_was_describe_bool
  \AddToHookNext{para/begin}{
    \bool_set_false:N \l_ptxtools_last_was_describe_bool
    \OmitIndent
  }
  \smallskip
}

% to support small layout adjustments for the texdoc version
\newcommand*{\IfOnlyDocumentation}[1]{
  \bool_if_exist:NT \g__codedoc_typeset_implementation_bool {
    \bool_if:NF \g__codedoc_typeset_implementation_bool {#1}
  }
}

\NewDocumentCommand{\PTXToolsDocSection}{om}{
  \str_if_eq:NNF  \CurrentFile \OTHRDocDTXfiles {
    \section[#1]{#2}
  }
}
\ExplSyntaxOff

\usepackage{underscore-ltx}

\DeclareRobustCommand{\issue}[1]{%
  {\DeclareFieldFormat{citeurlpostfix}{\href{##1/issues/#1}{\##1}}\citefield{othr-ci-repo}[citeurlpostfix]{url}}%
}
\newcommand*{\version}[1]{v#1}

\EnableCrossrefs
% \CodelineIndex
\RecordChanges

\DoNotIndex{\newcommand,\newenvironment}

\expandafter\ifx\csname OTHRDocDTXfiles\endcsname\relax
  \title{OTHRegensburg-CI --- \LaTeX\ using the Corporate Design of OTH-Regensburg%
    \thanks{This file describes version \fileversion, last revised \filedate.}
  }
  \GetFileInfo{othr-ci.dtx}
\else
  \GetFileInfo{\OTHRDocDTXfiles}
% subtitle? \fileinfo?
\fi

\author{%
  Marei Peischl\thanks{Email: \href{mailto:othr-ci@peitex.de}{othr-ci@peitex.de}}%
  \and Michael Niemetz\thanks{Email: \href{mailto:michael.niemetz@oth-regensburg.de}{michael.niemetz@oth-regensburg.de}}%
}
\date{\fileversion~from \filedate}

\usepackage{OTHR}

\providecommand*{\OTHRDocDTXfiles}{
  othr-ci.dtx,
  othr-beamer.dtx,
  othr-colors.dtx,
  othr-fonts.dtx
}

\begin{document}

\maketitle
\EnableDocumentation\DisableImplementation

\begin{abstract}
  The OTHRegensburg-CI bundle provides possibilities to use the CI of OTH Regensburg for \LaTeX{}.

  Parts of this bundle are historically grown, but the beamer template was completely rewritten in 2025. This document is the starting point for a full documentation of the bundle, but all document types beside beamer are described within the old documentation, which is only available in German \cite{othrTeX}.
\end{abstract}

\tableofcontents

\iffalse
\section{Contained packages and supported document types}

This bundle supports the following document types:

\begin{description}
\item[Beamer slides], see section \ref{sec:beamer}
\end{description}
\fi

\UseName{exp_args:No} \DocInput {\OTHRDocDTXfiles}

\DisableDocumentation\EnableImplementation

\begin{implementation}
\section{Implementation}
\end{implementation}

\DocInputAgain

\PrintChanges
\PrintIndex
\end{document}
%</driver>
% \fi
% \iffalse
%<*declareoptions|options>
%<*class>
% \fi
% \begin{implementation}
%    \begin{macrocode}
\str_const:Nn \c_@@_base_str {
%<beamer>  beamer
}

\exp_args:Ne \keys_define:nn {ptxcd/\c__ptxcd_base_str} {
%    \end{macrocode}
% \iffalse
%</class>
%    \begin{macrocode}
%<wrapper>  \keys_define:nn {ptxcd/othr-ci} {
%<gobble>  }^^A gobble brace for latexindent
%    \end{macrocode}
% \fi
% \end{implementation}
% \begin{documentation}
% \section{General layout options}
% The OTHRegensburg-CI bundle supports multiple document types.
% While rewriting the beamer theme the options which are available within multiple classes or packages of this bundle have been constructe separately to simplfy maintenance.
%
% Currently these options are only fully supporte by the beamer class \cls{othr-beamer} but this might be extended in the future.
% \end{documentation}
% \begin{documentation}
% \DescribeKeyOption{department=\meta{departmentshorthand}}{default}
% As the CI guideline is attaching colors to specific departments as well as there is a mechanism to load department specific config files this option may be used in a lot of different ways.
% For example the \pkg{othr-colors} package is using it to set  a department specific highlighting color.%^^A TODO provide reference!!
% The \value{default} is used to determine if a user specific setting as provided or not.
% In the current setup it has no effect towards leaving it empty.
% \end{documentation}
% \begin{implementation}
% \begin{optionenv}{department}
%    \begin{macrocode}
  department .choice:,
  department / default .code:n = \str_gset:Nn \g_ptxcd_department_str {default},
%<class>  department .initial:n = default,
  department / unknown .code:n = {
      \str_gset:NV \g_ptxcd_department_str \l_keys_value_tl
    },
%    \end{macrocode}
% \end{optionenv}
% \end{implementation}
% \begin{documentation}
% \DescribeKeyOption{departmentconfigprefix=\meta{prefix}}{othr-}
% The loading mechanism to be used to load department specific config files allows to change the file prefix.
% This might be useful if a specific department wants some internal templates to be shared separately.
% In that case it is also possible to use this option to set a specific directory where those config files can be found.
% By default this setting will search for files called |othr-|\meta{department}|.cfg|, where |othr-| can be set using this option and \option{department} would configure the second part.
% \end{documentation}
% \begin{implementation}
% \begin{optionenv}{departmentconfigprefix}
%    \begin{macrocode}
  departmentconfigprefix .tl_gset:N = \g_@@_config_prefix_tl,
%<class>  departmentconfigprefix .initial:n = othr-,
%    \end{macrocode}
% \end{optionenv}
% \end{implementation}
% \begin{documentation}
% \DescribeKeyOption{logofile=\meta{<main-logo-filename/prefix>}}{othr-logo}
% \DescribeKeyOption{titlelogofile=\meta{filename}}{\meta{logofile}_flat}
% \DescribeKeyOption{logofile=\meta{filename}}{\meta{logofile}_short}
% The beamer template is using two different logo files.
% The \enquote{flat} variant is used for the title page and the \enquote{short} appears within the frame title.
% To simplify the global config the option \option{logofile} is used as a wrapper to configure both defaults.
%
% As the official logos are shipped within this bundle it should not be necessary for users to use these options.
%
% \end{documentation}
% \begin{implementation}
% \begin{optionenv}{logofile,short-logofile,flat-logofile}
%    \begin{macrocode}
  logofile .tl_gset:N = \g_ptxcd_logofile_tl,
%<class>  logofile .initial:n = othr-logo,
  short-logofile .tl_gset:N = \g_@@_logofile_short_tl,
%<class>  short-logofile .initial:n =  \g_ptxcd_logofile_tl _short,
  flat-logofile .tl_gset:N = \g_@@_logofile_flat_tl,
%<class>  flat-logofile .initial:n =  \g_ptxcd_logofile_tl _flat,
%    \end{macrocode}
% \end{optionenv}
% \end{implementation}
% \begin{documentation}
% \DescribeKeyOption{accept-missing-logos=\meta{boolean}}{false}
% This option exists for compatibility reasons.
% It allows to use the \pkg{othr-logos} package which has been existing the former logo variants.
% In case you need to compile old tex-files and get error messages about missing logos, this option can be enabled.
% \end{documentation}
% \begin{implementation}
% \begin{optionenv}{accept-missing-logos}
%    \begin{macrocode}
  accept-missing-logos .bool_gset:N = \g_ptxcd_logo_workaround_bool,
%<class>  accept-missing-logos .initial:n = false,
  accept-missing-logos .usage:n = load,
  accept-missing-logos .default:n = true,
%    \end{macrocode}
% \end{optionenv}
% \iffalse
%<*wrapper>
% \fi
% \begin{optionenv}{useDepartmentLogo,useOTHRLogo}
%   Compatibility options for logo package.
%    \begin{macrocode}
  useDepartmentLogo .code:n = \PassOptionsToPackage{useDepartmentLogo}{othr-logos},
  useOTHRLogo .code:n = \PassOptionsToPackage{useOTHRLogo}{othr-logos},
%    \end{macrocode}
% \end{optionenv}
% \iffalse
%</wrapper>
%</declareoptions|options>
% \fi
%    \begin{macrocode}
%<processoptions|init-department>}
%    \end{macrocode}
% \iffalse
%<*init-department>
% \fi
% Configure defaults, in case the packages are used without loading a class being part of  OTHR-CI.
%    \begin{macrocode}
\str_if_exist:NF \g_ptxcd_department_str {
  \str_new:N \g_ptxcd_department_str
  \str_gset:Nn \g_ptxcd_department_str {default}
  \tl_gset:Nn \g_ptxcd_logofile_tl {othr-logo}
  \tl_gset:Nn \g_@@_logofile_short_tl { \g_ptxcd_logofile_tl _short}
  \tl_gset:Nn \g_@@_logofile_flat_tl { \g_ptxcd_logofile_tl _flat}
  \tl_gset:Nn \g_@@_config_prefix_tl {othr-}
}
%    \end{macrocode}
% \iffalse
%</init-department>
%<*processoptions>
% \fi
%    \begin{macrocode}
%<!wrapper>\ProcessKeyOptions[ptxcd/\c__ptxcd_base_str]
%<wrapper>\ProcessKeyOptions[ptxcd/othr-ci]
%    \end{macrocode}
% \iffalse
%</processoptions>
%<*wrapper&main>
% \fi
%    \begin{macrocode}
\RequirePackage{othr-colors}
\RequirePackage{othr-logos}
\RequirePackage{othr-fonts}
%    \end{macrocode}
% \iffalse
%</wrapper&main>
% \fi
% \end{implementation}
\endinput

%</driver>
% \fi
% \PTXToolsDocSection{Colors}
% \begin{documentation}
%
% The colors can be used like within standard \LaTeX.
% The predefined OTHR-CI specific colors are:
% \begin{multicols}{4}
%   \raggedright
%   \ExplSyntaxOn
%   \RenewDocumentCommand{\ProvidesFile}{mo}{}
%   \RenewDocumentCommand{\definecolor}{ommm}{\makebox[5em]{#2}\textcolor{#2}{\rule{1cm}{\ht\strutbox}}\par\ignorespaces}
%   \let\endinput\relax
%   \file_get:nnN {othr-colors.def} {\cctab_select:N \c_code_cctab} \l_tmpa_tl
%   \l_tmpa_tl
%   \ExplSyntaxOff
% \end{multicols}
% \end{documentation}
% \begin{implementation}
% \iffalse
%<*declareoptions>
% \fi
%    \begin{macrocode}
\str_if_exist:NF \g_ptxcd_department_str {
  \str_new:N \g_ptxcd_department_str
  \str_gset:Nn \g_ptxcd_department_str {default}
}

\bool_new:N \g_@@_color_blacktext_bool
\keys_define:nn {ptxcd/colors} {
%    \end{macrocode}
% \iffalse
%</declareoptions>
%<*options>
% \fi
% \end{implementation}
%
% \begin{documentation}
% \DescribeKeyOption{accentcolor=\meta{Color}}{0b}*
% \DescribeOption{accent=}
%
% Highlight color used for all highlight elements by default.
% The options below will reference this color as  default setting meta{accentcolor}.
% Depending on the documentclass it might be used for the so-called \enquote{identbar} as well as highlighted test elements, e.\,g. within letters if \option{premium} is enabled.
%
% |accent| is an alias.
% We recommend using |accentcolor| to simplify recognizing the purpose of that option.
% \end{documentation}
%
% \begin{implementation}
% \begin{optionenv}{accent=}
%    \begin{macrocode}
  accent .tl_gset:N = \g_ptxcd_color_accent_tl,
%<package>  accent .initial:n = 0b,
  accentcolor .meta:n = {accent =#1},
%    \end{macrocode}
% \end{optionenv}^^A accent
% \end{implementation}
%
% \begin{documentation}
% \DescribeKeyOption{textaccent=\meta{color}}{\meta{accentcolor}}*
% The accent color for highlighted text. Usage depends on the document class.
% \end{documentation}
%
% \begin{implementation}
% \begin{optionenv}{textaccent=\meta{color}}
%    \begin{macrocode}
  textaccent .tl_gset:N = \g_ptxcd_color_textaccent_tl,
%<package>  textaccent .initial:n = \g_ptxcd_color_accent_tl,
  textaccentcolor .meta:n = {textaccent = #1},
%    \end{macrocode}
% \end{optionenv}^^A textaccent
% \end{implementation}
%
% \begin{documentation}
% \DescribeKeyOption{text=black/default}{default}*
% The text placed on top a colored area, e.\,g. colored title background, is defined to have enough contrast.
% \end{documentation}
%
% \begin{implementation}
% \begin{optionenv}{text}
%   \iffalse
%<*package>
% \fi
%    \begin{macrocode}
  text .choice:,
  text/black .code:n = \bool_gset_true:N \g_@@_color_blacktext_bool,
  text/default .code:n = \bool_gset_false:N \g_@@_color_blacktext_bool,
  text .initial:n = {default},
%    \end{macrocode}
%   \iffalse
%</package>
% \fi
%    \begin{macrocode}
%<!package>  text .code:n = \exp_args:Ne \PassOptionsToPackage{text=\l_keys_value_tl}{othr-colors},
%    \end{macrocode}
%   For compatibility reasons there is an alias for \option{text=black} called \option{blackFont}.
%    \begin{macrocode}
  blackFont.meta:n = {text=black},
  blackFont .value_forbidden:n = true,
%    \end{macrocode}
% \end{optionenv}^^A text
% \end{implementation}
% \begin{documentation}
% \DescribeKeyOption*{colormode=cmyk/RGB}{Depends on document type}*
% The colormode can be selected to overwrite the default setting.
% Please be aware, that there is a difference in casing and \code{RGB} has to be used uppercase and \code{cmyk} lowercase.
% The reason for this is that internally uppercase and lowercase model names slightly differ because of rounding.
% \end{documentation}
% \begin{implementation}
% \begin{optionenv}{colormode}
%    \begin{macrocode}
  colormode .choices:nn = {cmyk,RGB}{
      \PassOptionsToPackage{\l_keys_value_tl}{xcolor}
    },
  colormode .usage:n = load,
  RGB .meta:n = {colormode=RGB},
  cymk .meta:n = {colormode=cmyk},
%    \end{macrocode}
% \end{optionenv}
% \iffalse
%<*package>
% \fi
%    \begin{macrocode}
}
\ProcessKeyOptions[ptxcd/colors]
\RequirePackage{xcolor}

% \iffalse meta-comment
%
% OTHR-CI -- LaTeX using the Corporate Design of OTH Regensburg
% ----------------------------------------------------------------------------
%
%  Copyright (C) 2015-2024 Michael Nietmetz <michael.niemetz@oth-regensburg.de>
%                 and the OTH Regensburg-LaTeX-Team
%                2025 Marei Peischl (peiTeX) <othr-ci@peitex.de>
%                 and Michael Nietmetz <michael.niemetz@oth-regensburg.de>
%
% ============================================================================
% This work may be distributed and/or modified under the
% conditions of the LaTeX Project Public License, either version 1.3c
% of this license or (at your option) any later version.
% The latest version of this license is in
% http://www.latex-project.org/lppl.txt
% and version 1.3c or later is part of all distributions of LaTeX
% version 2008/05/04 or later.
%
% This work has the LPPL maintenance status `maintained'.
%
% The Current Maintainers of this work are
%   Marei Peischl <othr-ci@peitex.de>
%   Michael Niemetz <michael.niemetz@oth-regensburg.de>
%
% The development repository can be found at
% https://codeberg.org/OTH-Regensburg/OTHRegensburg-LaTeX
% Please use the issue tracker for feedback!
% ============================================================================
%
% \fi
% \iffalse
%<*driver>
\ProvidesFile{othr-colors.dtx}
  [2025-05-22 v0.01-dev Color setup for  CI of OTH Regensburg]
%</driver>
%<@@=ptxcd>
%<*identification>
%<package>\NeedsTeXFormat{LaTeX2e}[2022-10-01]
%<package>\ProvidesExplPackage{othr-colors}{2025-05-22}{0.01-dev}{Color setup for  CI of OTH Regensburg}
%<definitions>\ProvidesFile{othr-colors.def}[2025-05-22 v0.01-dev Color definitions for OTHR-CI]
%</identification>
%<*driver>
\title{\textsf{othr-colors} --Color setup for the Corporate Design of OTH Regensburg%
  \thanks{This file describes version \fileversion, last revised \filedate.}
}
\providecommand*{\OTHRDocDTXfiles}{othr-colors.dtx}
\input{othr-ci.dtx}
%</driver>
% \fi
% \PTXToolsDocSection{Colors}
% \begin{documentation}
%
% The colors can be used like within standard \LaTeX.
% The predefined OTHR-CI specific colors are:
% \begin{multicols}{4}
%   \raggedright
%   \ExplSyntaxOn
%   \RenewDocumentCommand{\ProvidesFile}{mo}{}
%   \RenewDocumentCommand{\definecolor}{ommm}{\makebox[5em]{#2}\textcolor{#2}{\rule{1cm}{\ht\strutbox}}\par\ignorespaces}
%   \let\endinput\relax
%   \file_get:nnN {othr-colors.def} {\cctab_select:N \c_code_cctab} \l_tmpa_tl
%   \l_tmpa_tl
%   \ExplSyntaxOff
% \end{multicols}
% \end{documentation}
% \begin{implementation}
% \iffalse
%<*declareoptions>
% \fi
%    \begin{macrocode}
\str_if_exist:NF \g_ptxcd_department_str {
  \str_new:N \g_ptxcd_department_str
  \str_gset:Nn \g_ptxcd_department_str {default}
}

\bool_new:N \g_@@_color_blacktext_bool
\keys_define:nn {ptxcd/colors} {
%    \end{macrocode}
% \iffalse
%</declareoptions>
%<*options>
% \fi
% \end{implementation}
%
% \begin{documentation}
% \DescribeKeyOption{accentcolor=\meta{Color}}{0b}*
% \DescribeOption{accent=}
%
% Highlight color used for all highlight elements by default.
% The options below will reference this color as  default setting meta{accentcolor}.
% Depending on the documentclass it might be used for the so-called \enquote{identbar} as well as highlighted test elements, e.\,g. within letters if \option{premium} is enabled.
%
% |accent| is an alias.
% We recommend using |accentcolor| to simplify recognizing the purpose of that option.
% \end{documentation}
%
% \begin{implementation}
% \begin{optionenv}{accent=}
%    \begin{macrocode}
  accent .tl_gset:N = \g_ptxcd_color_accent_tl,
%<package>  accent .initial:n = 0b,
  accentcolor .meta:n = {accent =#1},
%    \end{macrocode}
% \end{optionenv}^^A accent
% \end{implementation}
%
% \begin{documentation}
% \DescribeKeyOption{textaccent=\meta{color}}{\meta{accentcolor}}*
% The accent color for highlighted text. Usage depends on the document class.
% \end{documentation}
%
% \begin{implementation}
% \begin{optionenv}{textaccent=\meta{color}}
%    \begin{macrocode}
  textaccent .tl_gset:N = \g_ptxcd_color_textaccent_tl,
%<package>  textaccent .initial:n = \g_ptxcd_color_accent_tl,
  textaccentcolor .meta:n = {textaccent = #1},
%    \end{macrocode}
% \end{optionenv}^^A textaccent
% \end{implementation}
%
% \begin{documentation}
% \DescribeKeyOption{text=black/default}{default}*
% The text placed on top a colored area, e.\,g. colored title background, is defined to have enough contrast.
% \end{documentation}
%
% \begin{implementation}
% \begin{optionenv}{text}
%   \iffalse
%<*package>
% \fi
%    \begin{macrocode}
  text .choice:,
  text/black .code:n = \bool_gset_true:N \g_@@_color_blacktext_bool,
  text/default .code:n = \bool_gset_false:N \g_@@_color_blacktext_bool,
  text .initial:n = {default},
%    \end{macrocode}
%   \iffalse
%</package>
% \fi
%    \begin{macrocode}
%<!package>  text .code:n = \exp_args:Ne \PassOptionsToPackage{text=\l_keys_value_tl}{othr-colors},
%    \end{macrocode}
%   For compatibility reasons there is an alias for \option{text=black} called \option{blackFont}.
%    \begin{macrocode}
  blackFont.meta:n = {text=black},
  blackFont .value_forbidden:n = true,
%    \end{macrocode}
% \end{optionenv}^^A text
% \end{implementation}
% \begin{documentation}
% \DescribeKeyOption*{colormode=cmyk/RGB}{Depends on document type}*
% The colormode can be selected to overwrite the default setting.
% Please be aware, that there is a difference in casing and \code{RGB} has to be used uppercase and \code{cmyk} lowercase.
% The reason for this is that internally uppercase and lowercase model names slightly differ because of rounding.
% \end{documentation}
% \begin{implementation}
% \begin{optionenv}{colormode}
%    \begin{macrocode}
  colormode .choices:nn = {cmyk,RGB}{
      \PassOptionsToPackage{\l_keys_value_tl}{xcolor}
    },
  colormode .usage:n = load,
  RGB .meta:n = {colormode=RGB},
  cymk .meta:n = {colormode=cmyk},
%    \end{macrocode}
% \end{optionenv}
% \iffalse
%<*package>
% \fi
%    \begin{macrocode}
}
\ProcessKeyOptions[ptxcd/colors]
\RequirePackage{xcolor}

\input{othr-colors.def}
%    \end{macrocode}
% \iffalse
%</package>
%</options>
%<*definitions>
% \fi
%    \begin{macrocode}

\definecolor{OTHRgrau}{cmyk/RGB/HTML}{0,0,0,.5/157,157,156/9D9D9C}

\definecolor{FakANK}{cmyk/RGB/HTML}{.03,.03,.15,.20/213,210,194/D5D2C2}
\definecolor{FakA}{cmyk/RGB/HTML}{0,0,1,0/255,237,0/FFED00}
\definecolor{FakB}{cmyk/RGB/HTML}{0,.6,1,0/239,125,0/EF7D00}
\definecolor{FakBM}{cmyk/RGB/HTML}{0,.9,1,.2/194,46,12/C22E0C}
\definecolor{FakEI}{cmyk/RGB/HTML}{.6,1,0,0/131,31,130/831F82}
\definecolor{FakIM}{cmyk/RGB/HTML}{1,.8,0,0/22,65,148/164194}
\definecolor{FakM}{cmyk/RGB/HTML}{1,0,.35,0/0,155,172/009BAC}
\definecolor{FakSG}{cmyk/RGB/HTML}{1,0,.9,.1/0,140,73/008C49}

% Fakultäten (Grundfarben), die nicht im Handbuch vorkommen
\definecolor{FakAM}{cmyk}{.03,.03,.15,.2}
\definecolor{FakBW}{cmyk}{0,.9,1,.2}
\definecolor{FakS}{cmyk}{1,0,.9,.1}

\definecolor{ZWW}{cmyk}{.2,0,1,0}
\definecolor{LASIII}{rgb}{.223,.412,.489}
\definecolor{LASIIIcontrast}{rgb}{.737,.850,.482}

%    \end{macrocode}
% \iffalse
%</definitions>
%<*package&body>
% \fi
%    \begin{macrocode}

\ExplSyntaxOn
% 1. Arg: Basisname der Farbe
% 2. Arg: Farbe von der abgeleitet wird
% 3. Arg: Schriftfarbe
% 4. Arg: Kontrastfarbe
\newcommand{\@createColors}[4]{
  \colorlet{#1}{#2}
  \colorlet{#125}{#1!25}
  \colorlet{#150}{#1!50}
  \colorlet{#190}{#1!90}
  \bool_if:NTF \g_@@_color_blacktext_bool {
% wenn die Kontrastfarbe nicht weiss ist,
% so darf sie auch bei blackFont nicht weiss sein
    \str_if_eq:nnTF {#4} {white} {
      \colorlet{#1FontColor}{black}
      \colorlet{#1ContrastColor}{white}
    }{
      \colorlet{#1FontColor}{white}
      \colorlet{#1ContrastColor}{black}
    }
  }{
    \colorlet{#1FontColor}{#3}
    \colorlet{#1ContrastColor}{#4}
  }
}
\ExplSyntaxOff

\newcommand{\@createCurrentDepartmentColors}[1]{
  \colorlet{Accent}{#1}
  \colorlet{Accent25}{#125}
  \colorlet{Accent50}{#150}
  \colorlet{Accent90}{#190}
  \colorlet{AccentFontColor}{#1FontColor}
  \colorlet{AccentContrastColor}{#1ContrastColor}
}

% OTHR Farben
\@createColors{default}{OTHRgrau}{OTHRgrau}{white}
\@createColors{OTHR}{OTHRgrau}{OTHRgrau}{white}

% sonstiges
\@createColors{ZWW}{ZWW}{OTHRgrau}{black}

% Fakultätsfarben
\@createColors{FakA}{FakA}{OTHRgrau}{black}
\@createColors{FakAM}{FakAM}{FakAM}{black}
\@createColors{FakANK}{FakANK}{FakANK}{black}
\@createColors{FakB}{FakB}{FakB}{white}
\@createColors{FakBW}{FakBW}{FakBW}{white}
\@createColors{FakBM}{FakBM}{FakBM}{white}
\@createColors{FakEI}{FakEI}{FakEI}{white}
\@createColors{FakIM}{FakIM}{FakIM}{white}
\@createColors{FakM}{FakM}{FakM}{white}
\@createColors{FakS}{FakS}{FakS}{white}
\@createColors{FakSG}{FakSG}{FakSG}{white}

% % Centers and Schools
\@createColors{RCER}{OTHRgrau}{OTHRgrau}{white}
\@createColors{RCHST}{OTHRgrau}{OTHRgrau}{white}
\@createColors{RCAI}{OTHRgrau}{OTHRgrau}{white}
\@createColors{RSDS}{OTHRgrau}{OTHRgrau}{white}

% Laborfarben
% FakA
\@createColors{FMIS}{FakA}{OTHRgrau}{black}

% FakAM
\@createColors{NACH}{FakANK}{FakANK}{white}
\@createColors{SappZ}{FakANK}{FakANK}{white}

% FakB
\@createColors{Geo}{FakB}{FakB}{white}
\@createColors{KNB}{FakB}{FakB}{white}

% Fak BW
\@createColors{12Science}{FakBW}{FakBW}{white}
\@createColors{FEURO}{FakBW}{FakBW}{white}

% FakEI
\@createColors{Bisp}{FakEI}{FakEI}{white}
\@createColors{LAS3}{LASIII}{LASIII}{LASIIIcontrast}
\@createColors{DK0PT}{FakEI}{FakEI}{white}
\@createColors{FENES}{FakEI}{FakEI}{white}
\@createColors{MRU}{FakEI}{FakEI}{white}
\@createColors{SES}{FakEI}{FakEI}{white}

% FakIM
\@createColors{CCSE}{FakIM}{FakIM}{white}
\@createColors{eHealth}{FakIM}{FakIM}{white}
\@createColors{ITZ}{FakIM}{FakIM}{white}
\@createColors{LFD}{FakIM}{FakIM}{white}
\@createColors{MD}{FakIM}{FakIM}{white}
\@createColors{ReMIC}{FakIM}{FakIM}{white}
\@createColors{SEC}{FakIM}{FakIM}{white}
\@createColors{IMDS}{FakIM}{FakIM}{white}

% FakM
\@createColors{AIMS}{FakM}{FakM}{white}
\@createColors{BFM}{FakM}{FakM}{white}
\@createColors{BMA}{FakM}{FakM}{white}
\@createColors{CEEC}{FakM}{FakM}{white}
\@createColors{CFD}{FakM}{FakM}{white}
\@createColors{IPF}{FakM}{FakM}{white}
\@createColors{KIB}{FakM}{FakM}{white}
\@createColors{KWK}{FakM}{FakM}{white}
\@createColors{LAT}{FakM}{FakM}{white}
\@createColors{LBM}{FakM}{FakM}{white}
\@createColors{LeanLab}{FakM}{FakM}{white}
\@createColors{LFT}{FakM}{FakM}{white}
\@createColors{LFW}{FakM}{FakM}{white}
\@createColors{LMP}{FakM}{FakM}{white}
\@createColors{LMS}{FakM}{FakM}{white}
\@createColors{LRT}{FakM}{FakM}{white}
\@createColors{LWM}{FakM}{FakM}{white}
\@createColors{LWS}{FakM}{FakM}{white}
\@createColors{MKS}{FakM}{FakM}{white}
\@createColors{MST}{FakM}{FakM}{white}
\@createColors{RRRU}{FakM}{FakM}{white}
\@createColors{RST}{FakM}{FakM}{white}
\@createColors{FEM}{FakM}{FakM}{white}

% FakS
\@createColors{IST}{FakS}{FakS}{white}
\@createColors{LP}{FakS}{FakS}{white}
\@createColors{PT}{FakS}{FakS}{white}
\@createColors{PF}{FakS}{FakS}{white}

% current Department
\@createCurrentDepartmentColors{\UseName{g_ptxcd_department_str}}

%    \end{macrocode}
% \iffalse
%</package&body>
% \fi
% \end{implementation}
\endinput

%    \end{macrocode}
% \iffalse
%</package>
%</options>
%<*definitions>
% \fi
%    \begin{macrocode}

\definecolor{OTHRgrau}{cmyk/RGB/HTML}{0,0,0,.5/157,157,156/9D9D9C}

\definecolor{FakANK}{cmyk/RGB/HTML}{.03,.03,.15,.20/213,210,194/D5D2C2}
\definecolor{FakA}{cmyk/RGB/HTML}{0,0,1,0/255,237,0/FFED00}
\definecolor{FakB}{cmyk/RGB/HTML}{0,.6,1,0/239,125,0/EF7D00}
\definecolor{FakBM}{cmyk/RGB/HTML}{0,.9,1,.2/194,46,12/C22E0C}
\definecolor{FakEI}{cmyk/RGB/HTML}{.6,1,0,0/131,31,130/831F82}
\definecolor{FakIM}{cmyk/RGB/HTML}{1,.8,0,0/22,65,148/164194}
\definecolor{FakM}{cmyk/RGB/HTML}{1,0,.35,0/0,155,172/009BAC}
\definecolor{FakSG}{cmyk/RGB/HTML}{1,0,.9,.1/0,140,73/008C49}

% Fakultäten (Grundfarben), die nicht im Handbuch vorkommen
\definecolor{FakAM}{cmyk}{.03,.03,.15,.2}
\definecolor{FakBW}{cmyk}{0,.9,1,.2}
\definecolor{FakS}{cmyk}{1,0,.9,.1}

\definecolor{ZWW}{cmyk}{.2,0,1,0}
\definecolor{LASIII}{rgb}{.223,.412,.489}
\definecolor{LASIIIcontrast}{rgb}{.737,.850,.482}

%    \end{macrocode}
% \iffalse
%</definitions>
%<*package&body>
% \fi
%    \begin{macrocode}

\ExplSyntaxOn
% 1. Arg: Basisname der Farbe
% 2. Arg: Farbe von der abgeleitet wird
% 3. Arg: Schriftfarbe
% 4. Arg: Kontrastfarbe
\newcommand{\@createColors}[4]{
  \colorlet{#1}{#2}
  \colorlet{#125}{#1!25}
  \colorlet{#150}{#1!50}
  \colorlet{#190}{#1!90}
  \bool_if:NTF \g_@@_color_blacktext_bool {
% wenn die Kontrastfarbe nicht weiss ist,
% so darf sie auch bei blackFont nicht weiss sein
    \str_if_eq:nnTF {#4} {white} {
      \colorlet{#1FontColor}{black}
      \colorlet{#1ContrastColor}{white}
    }{
      \colorlet{#1FontColor}{white}
      \colorlet{#1ContrastColor}{black}
    }
  }{
    \colorlet{#1FontColor}{#3}
    \colorlet{#1ContrastColor}{#4}
  }
}
\ExplSyntaxOff

\newcommand{\@createCurrentDepartmentColors}[1]{
  \colorlet{Accent}{#1}
  \colorlet{Accent25}{#125}
  \colorlet{Accent50}{#150}
  \colorlet{Accent90}{#190}
  \colorlet{AccentFontColor}{#1FontColor}
  \colorlet{AccentContrastColor}{#1ContrastColor}
}

% OTHR Farben
\@createColors{default}{OTHRgrau}{OTHRgrau}{white}
\@createColors{OTHR}{OTHRgrau}{OTHRgrau}{white}

% sonstiges
\@createColors{ZWW}{ZWW}{OTHRgrau}{black}

% Fakultätsfarben
\@createColors{FakA}{FakA}{OTHRgrau}{black}
\@createColors{FakAM}{FakAM}{FakAM}{black}
\@createColors{FakANK}{FakANK}{FakANK}{black}
\@createColors{FakB}{FakB}{FakB}{white}
\@createColors{FakBW}{FakBW}{FakBW}{white}
\@createColors{FakBM}{FakBM}{FakBM}{white}
\@createColors{FakEI}{FakEI}{FakEI}{white}
\@createColors{FakIM}{FakIM}{FakIM}{white}
\@createColors{FakM}{FakM}{FakM}{white}
\@createColors{FakS}{FakS}{FakS}{white}
\@createColors{FakSG}{FakSG}{FakSG}{white}

% % Centers and Schools
\@createColors{RCER}{OTHRgrau}{OTHRgrau}{white}
\@createColors{RCHST}{OTHRgrau}{OTHRgrau}{white}
\@createColors{RCAI}{OTHRgrau}{OTHRgrau}{white}
\@createColors{RSDS}{OTHRgrau}{OTHRgrau}{white}

% Laborfarben
% FakA
\@createColors{FMIS}{FakA}{OTHRgrau}{black}

% FakAM
\@createColors{NACH}{FakANK}{FakANK}{white}
\@createColors{SappZ}{FakANK}{FakANK}{white}

% FakB
\@createColors{Geo}{FakB}{FakB}{white}
\@createColors{KNB}{FakB}{FakB}{white}

% Fak BW
\@createColors{12Science}{FakBW}{FakBW}{white}
\@createColors{FEURO}{FakBW}{FakBW}{white}

% FakEI
\@createColors{Bisp}{FakEI}{FakEI}{white}
\@createColors{LAS3}{LASIII}{LASIII}{LASIIIcontrast}
\@createColors{DK0PT}{FakEI}{FakEI}{white}
\@createColors{FENES}{FakEI}{FakEI}{white}
\@createColors{MRU}{FakEI}{FakEI}{white}
\@createColors{SES}{FakEI}{FakEI}{white}

% FakIM
\@createColors{CCSE}{FakIM}{FakIM}{white}
\@createColors{eHealth}{FakIM}{FakIM}{white}
\@createColors{ITZ}{FakIM}{FakIM}{white}
\@createColors{LFD}{FakIM}{FakIM}{white}
\@createColors{MD}{FakIM}{FakIM}{white}
\@createColors{ReMIC}{FakIM}{FakIM}{white}
\@createColors{SEC}{FakIM}{FakIM}{white}
\@createColors{IMDS}{FakIM}{FakIM}{white}

% FakM
\@createColors{AIMS}{FakM}{FakM}{white}
\@createColors{BFM}{FakM}{FakM}{white}
\@createColors{BMA}{FakM}{FakM}{white}
\@createColors{CEEC}{FakM}{FakM}{white}
\@createColors{CFD}{FakM}{FakM}{white}
\@createColors{IPF}{FakM}{FakM}{white}
\@createColors{KIB}{FakM}{FakM}{white}
\@createColors{KWK}{FakM}{FakM}{white}
\@createColors{LAT}{FakM}{FakM}{white}
\@createColors{LBM}{FakM}{FakM}{white}
\@createColors{LeanLab}{FakM}{FakM}{white}
\@createColors{LFT}{FakM}{FakM}{white}
\@createColors{LFW}{FakM}{FakM}{white}
\@createColors{LMP}{FakM}{FakM}{white}
\@createColors{LMS}{FakM}{FakM}{white}
\@createColors{LRT}{FakM}{FakM}{white}
\@createColors{LWM}{FakM}{FakM}{white}
\@createColors{LWS}{FakM}{FakM}{white}
\@createColors{MKS}{FakM}{FakM}{white}
\@createColors{MST}{FakM}{FakM}{white}
\@createColors{RRRU}{FakM}{FakM}{white}
\@createColors{RST}{FakM}{FakM}{white}
\@createColors{FEM}{FakM}{FakM}{white}

% FakS
\@createColors{IST}{FakS}{FakS}{white}
\@createColors{LP}{FakS}{FakS}{white}
\@createColors{PT}{FakS}{FakS}{white}
\@createColors{PF}{FakS}{FakS}{white}

% current Department
\@createCurrentDepartmentColors{\UseName{g_ptxcd_department_str}}

%    \end{macrocode}
% \iffalse
%</package&body>
% \fi
% \end{implementation}
\endinput

%    \end{macrocode}
% \iffalse
%</package>
%</options>
%<*definitions>
% \fi
%    \begin{macrocode}

\definecolor{OTHRgrau}{cmyk/RGB/HTML}{0,0,0,.5/157,157,156/9D9D9C}

\definecolor{FakANK}{cmyk/RGB/HTML}{.03,.03,.15,.20/213,210,194/D5D2C2}
\definecolor{FakA}{cmyk/RGB/HTML}{0,0,1,0/255,237,0/FFED00}
\definecolor{FakB}{cmyk/RGB/HTML}{0,.6,1,0/239,125,0/EF7D00}
\definecolor{FakBM}{cmyk/RGB/HTML}{0,.9,1,.2/194,46,12/C22E0C}
\definecolor{FakEI}{cmyk/RGB/HTML}{.6,1,0,0/131,31,130/831F82}
\definecolor{FakIM}{cmyk/RGB/HTML}{1,.8,0,0/22,65,148/164194}
\definecolor{FakM}{cmyk/RGB/HTML}{1,0,.35,0/0,155,172/009BAC}
\definecolor{FakSG}{cmyk/RGB/HTML}{1,0,.9,.1/0,140,73/008C49}

% Fakultäten (Grundfarben), die nicht im Handbuch vorkommen
\definecolor{FakAM}{cmyk}{.03,.03,.15,.2}
\definecolor{FakBW}{cmyk}{0,.9,1,.2}
\definecolor{FakS}{cmyk}{1,0,.9,.1}

\definecolor{ZWW}{cmyk}{.2,0,1,0}
\definecolor{LASIII}{rgb}{.223,.412,.489}
\definecolor{LASIIIcontrast}{rgb}{.737,.850,.482}

%    \end{macrocode}
% \iffalse
%</definitions>
%<*package&body>
% \fi
%    \begin{macrocode}

\ExplSyntaxOn
% 1. Arg: Basisname der Farbe
% 2. Arg: Farbe von der abgeleitet wird
% 3. Arg: Schriftfarbe
% 4. Arg: Kontrastfarbe
\newcommand{\@createColors}[4]{
  \colorlet{#1}{#2}
  \colorlet{#125}{#1!25}
  \colorlet{#150}{#1!50}
  \colorlet{#190}{#1!90}
  \bool_if:NTF \g_@@_color_blacktext_bool {
% wenn die Kontrastfarbe nicht weiss ist,
% so darf sie auch bei blackFont nicht weiss sein
    \str_if_eq:nnTF {#4} {white} {
      \colorlet{#1FontColor}{black}
      \colorlet{#1ContrastColor}{white}
    }{
      \colorlet{#1FontColor}{white}
      \colorlet{#1ContrastColor}{black}
    }
  }{
    \colorlet{#1FontColor}{#3}
    \colorlet{#1ContrastColor}{#4}
  }
}
\ExplSyntaxOff

\newcommand{\@createCurrentDepartmentColors}[1]{
  \colorlet{Accent}{#1}
  \colorlet{Accent25}{#125}
  \colorlet{Accent50}{#150}
  \colorlet{Accent90}{#190}
  \colorlet{AccentFontColor}{#1FontColor}
  \colorlet{AccentContrastColor}{#1ContrastColor}
}

% OTHR Farben
\@createColors{default}{OTHRgrau}{OTHRgrau}{white}
\@createColors{OTHR}{OTHRgrau}{OTHRgrau}{white}

% sonstiges
\@createColors{ZWW}{ZWW}{OTHRgrau}{black}

% Fakultätsfarben
\@createColors{FakA}{FakA}{OTHRgrau}{black}
\@createColors{FakAM}{FakAM}{FakAM}{black}
\@createColors{FakANK}{FakANK}{FakANK}{black}
\@createColors{FakB}{FakB}{FakB}{white}
\@createColors{FakBW}{FakBW}{FakBW}{white}
\@createColors{FakBM}{FakBM}{FakBM}{white}
\@createColors{FakEI}{FakEI}{FakEI}{white}
\@createColors{FakIM}{FakIM}{FakIM}{white}
\@createColors{FakM}{FakM}{FakM}{white}
\@createColors{FakS}{FakS}{FakS}{white}
\@createColors{FakSG}{FakSG}{FakSG}{white}

% % Centers and Schools
\@createColors{RCER}{OTHRgrau}{OTHRgrau}{white}
\@createColors{RCHST}{OTHRgrau}{OTHRgrau}{white}
\@createColors{RCAI}{OTHRgrau}{OTHRgrau}{white}
\@createColors{RSDS}{OTHRgrau}{OTHRgrau}{white}

% Laborfarben
% FakA
\@createColors{FMIS}{FakA}{OTHRgrau}{black}

% FakAM
\@createColors{NACH}{FakANK}{FakANK}{white}
\@createColors{SappZ}{FakANK}{FakANK}{white}

% FakB
\@createColors{Geo}{FakB}{FakB}{white}
\@createColors{KNB}{FakB}{FakB}{white}

% Fak BW
\@createColors{12Science}{FakBW}{FakBW}{white}
\@createColors{FEURO}{FakBW}{FakBW}{white}

% FakEI
\@createColors{Bisp}{FakEI}{FakEI}{white}
\@createColors{LAS3}{LASIII}{LASIII}{LASIIIcontrast}
\@createColors{DK0PT}{FakEI}{FakEI}{white}
\@createColors{FENES}{FakEI}{FakEI}{white}
\@createColors{MRU}{FakEI}{FakEI}{white}
\@createColors{SES}{FakEI}{FakEI}{white}

% FakIM
\@createColors{CCSE}{FakIM}{FakIM}{white}
\@createColors{eHealth}{FakIM}{FakIM}{white}
\@createColors{ITZ}{FakIM}{FakIM}{white}
\@createColors{LFD}{FakIM}{FakIM}{white}
\@createColors{MD}{FakIM}{FakIM}{white}
\@createColors{ReMIC}{FakIM}{FakIM}{white}
\@createColors{SEC}{FakIM}{FakIM}{white}
\@createColors{IMDS}{FakIM}{FakIM}{white}

% FakM
\@createColors{AIMS}{FakM}{FakM}{white}
\@createColors{BFM}{FakM}{FakM}{white}
\@createColors{BMA}{FakM}{FakM}{white}
\@createColors{CEEC}{FakM}{FakM}{white}
\@createColors{CFD}{FakM}{FakM}{white}
\@createColors{IPF}{FakM}{FakM}{white}
\@createColors{KIB}{FakM}{FakM}{white}
\@createColors{KWK}{FakM}{FakM}{white}
\@createColors{LAT}{FakM}{FakM}{white}
\@createColors{LBM}{FakM}{FakM}{white}
\@createColors{LeanLab}{FakM}{FakM}{white}
\@createColors{LFT}{FakM}{FakM}{white}
\@createColors{LFW}{FakM}{FakM}{white}
\@createColors{LMP}{FakM}{FakM}{white}
\@createColors{LMS}{FakM}{FakM}{white}
\@createColors{LRT}{FakM}{FakM}{white}
\@createColors{LWM}{FakM}{FakM}{white}
\@createColors{LWS}{FakM}{FakM}{white}
\@createColors{MKS}{FakM}{FakM}{white}
\@createColors{MST}{FakM}{FakM}{white}
\@createColors{RRRU}{FakM}{FakM}{white}
\@createColors{RST}{FakM}{FakM}{white}
\@createColors{FEM}{FakM}{FakM}{white}

% FakS
\@createColors{IST}{FakS}{FakS}{white}
\@createColors{LP}{FakS}{FakS}{white}
\@createColors{PT}{FakS}{FakS}{white}
\@createColors{PF}{FakS}{FakS}{white}

% current Department
\@createCurrentDepartmentColors{\UseName{g_ptxcd_department_str}}

%    \end{macrocode}
% \iffalse
%</package&body>
% \fi
% \end{implementation}
\endinput

%    \end{macrocode}
% \iffalse
%</package>
%</options>
%<*definitions>
% \fi
%    \begin{macrocode}

\definecolor{OTHRgrau}{cmyk/RGB/HTML}{0,0,0,.5/157,157,156/9D9D9C}

\definecolor{FakANK}{cmyk/RGB/HTML}{.03,.03,.15,.20/213,210,194/D5D2C2}
\definecolor{FakA}{cmyk/RGB/HTML}{0,0,1,0/255,237,0/FFED00}
\definecolor{FakB}{cmyk/RGB/HTML}{0,.6,1,0/239,125,0/EF7D00}
\definecolor{FakBM}{cmyk/RGB/HTML}{0,.9,1,.2/194,46,12/C22E0C}
\definecolor{FakEI}{cmyk/RGB/HTML}{.6,1,0,0/131,31,130/831F82}
\definecolor{FakIM}{cmyk/RGB/HTML}{1,.8,0,0/22,65,148/164194}
\definecolor{FakM}{cmyk/RGB/HTML}{1,0,.35,0/0,155,172/009BAC}
\definecolor{FakSG}{cmyk/RGB/HTML}{1,0,.9,.1/0,140,73/008C49}

% Fakultäten (Grundfarben), die nicht im Handbuch vorkommen
\definecolor{FakAM}{cmyk}{.03,.03,.15,.2}
\definecolor{FakBW}{cmyk}{0,.9,1,.2}
\definecolor{FakS}{cmyk}{1,0,.9,.1}

\definecolor{ZWW}{cmyk}{.2,0,1,0}
\definecolor{LASIII}{rgb}{.223,.412,.489}
\definecolor{LASIIIcontrast}{rgb}{.737,.850,.482}

%    \end{macrocode}
% \iffalse
%</definitions>
%<*package&body>
% \fi
%    \begin{macrocode}

\ExplSyntaxOn
% 1. Arg: Basisname der Farbe
% 2. Arg: Farbe von der abgeleitet wird
% 3. Arg: Schriftfarbe
% 4. Arg: Kontrastfarbe
\newcommand{\@createColors}[4]{
  \colorlet{#1}{#2}
  \colorlet{#125}{#1!25}
  \colorlet{#150}{#1!50}
  \colorlet{#190}{#1!90}
  \bool_if:NTF \g_@@_color_blacktext_bool {
% wenn die Kontrastfarbe nicht weiss ist,
% so darf sie auch bei blackFont nicht weiss sein
    \str_if_eq:nnTF {#4} {white} {
      \colorlet{#1FontColor}{black}
      \colorlet{#1ContrastColor}{white}
    }{
      \colorlet{#1FontColor}{white}
      \colorlet{#1ContrastColor}{black}
    }
  }{
    \colorlet{#1FontColor}{#3}
    \colorlet{#1ContrastColor}{#4}
  }
}
\ExplSyntaxOff

\newcommand{\@createCurrentDepartmentColors}[1]{
  \colorlet{Accent}{#1}
  \colorlet{Accent25}{#125}
  \colorlet{Accent50}{#150}
  \colorlet{Accent90}{#190}
  \colorlet{AccentFontColor}{#1FontColor}
  \colorlet{AccentContrastColor}{#1ContrastColor}
}

% OTHR Farben
\@createColors{default}{OTHRgrau}{OTHRgrau}{white}
\@createColors{OTHR}{OTHRgrau}{OTHRgrau}{white}

% sonstiges
\@createColors{ZWW}{ZWW}{OTHRgrau}{black}

% Fakultätsfarben
\@createColors{FakA}{FakA}{OTHRgrau}{black}
\@createColors{FakAM}{FakAM}{FakAM}{black}
\@createColors{FakANK}{FakANK}{FakANK}{black}
\@createColors{FakB}{FakB}{FakB}{white}
\@createColors{FakBW}{FakBW}{FakBW}{white}
\@createColors{FakBM}{FakBM}{FakBM}{white}
\@createColors{FakEI}{FakEI}{FakEI}{white}
\@createColors{FakIM}{FakIM}{FakIM}{white}
\@createColors{FakM}{FakM}{FakM}{white}
\@createColors{FakS}{FakS}{FakS}{white}
\@createColors{FakSG}{FakSG}{FakSG}{white}

% % Centers and Schools
\@createColors{RCER}{OTHRgrau}{OTHRgrau}{white}
\@createColors{RCHST}{OTHRgrau}{OTHRgrau}{white}
\@createColors{RCAI}{OTHRgrau}{OTHRgrau}{white}
\@createColors{RSDS}{OTHRgrau}{OTHRgrau}{white}

% Laborfarben
% FakA
\@createColors{FMIS}{FakA}{OTHRgrau}{black}

% FakAM
\@createColors{NACH}{FakANK}{FakANK}{white}
\@createColors{SappZ}{FakANK}{FakANK}{white}

% FakB
\@createColors{Geo}{FakB}{FakB}{white}
\@createColors{KNB}{FakB}{FakB}{white}

% Fak BW
\@createColors{12Science}{FakBW}{FakBW}{white}
\@createColors{FEURO}{FakBW}{FakBW}{white}

% FakEI
\@createColors{Bisp}{FakEI}{FakEI}{white}
\@createColors{LAS3}{LASIII}{LASIII}{LASIIIcontrast}
\@createColors{DK0PT}{FakEI}{FakEI}{white}
\@createColors{FENES}{FakEI}{FakEI}{white}
\@createColors{MRU}{FakEI}{FakEI}{white}
\@createColors{SES}{FakEI}{FakEI}{white}

% FakIM
\@createColors{CCSE}{FakIM}{FakIM}{white}
\@createColors{eHealth}{FakIM}{FakIM}{white}
\@createColors{ITZ}{FakIM}{FakIM}{white}
\@createColors{LFD}{FakIM}{FakIM}{white}
\@createColors{MD}{FakIM}{FakIM}{white}
\@createColors{ReMIC}{FakIM}{FakIM}{white}
\@createColors{SEC}{FakIM}{FakIM}{white}
\@createColors{IMDS}{FakIM}{FakIM}{white}

% FakM
\@createColors{AIMS}{FakM}{FakM}{white}
\@createColors{BFM}{FakM}{FakM}{white}
\@createColors{BMA}{FakM}{FakM}{white}
\@createColors{CEEC}{FakM}{FakM}{white}
\@createColors{CFD}{FakM}{FakM}{white}
\@createColors{IPF}{FakM}{FakM}{white}
\@createColors{KIB}{FakM}{FakM}{white}
\@createColors{KWK}{FakM}{FakM}{white}
\@createColors{LAT}{FakM}{FakM}{white}
\@createColors{LBM}{FakM}{FakM}{white}
\@createColors{LeanLab}{FakM}{FakM}{white}
\@createColors{LFT}{FakM}{FakM}{white}
\@createColors{LFW}{FakM}{FakM}{white}
\@createColors{LMP}{FakM}{FakM}{white}
\@createColors{LMS}{FakM}{FakM}{white}
\@createColors{LRT}{FakM}{FakM}{white}
\@createColors{LWM}{FakM}{FakM}{white}
\@createColors{LWS}{FakM}{FakM}{white}
\@createColors{MKS}{FakM}{FakM}{white}
\@createColors{MST}{FakM}{FakM}{white}
\@createColors{RRRU}{FakM}{FakM}{white}
\@createColors{RST}{FakM}{FakM}{white}
\@createColors{FEM}{FakM}{FakM}{white}

% FakS
\@createColors{IST}{FakS}{FakS}{white}
\@createColors{LP}{FakS}{FakS}{white}
\@createColors{PT}{FakS}{FakS}{white}
\@createColors{PF}{FakS}{FakS}{white}

% current Department
\@createCurrentDepartmentColors{\UseName{g_ptxcd_department_str}}

%    \end{macrocode}
% \iffalse
%</package&body>
% \fi
% \end{implementation}
\endinput
